% Source: Weinberg GR, physical PDF pp. 401--407 (printed pp. 381--385).
% Physical pp. 404--405 repeat physical pp. 402--403 (printed pp. 382--383);
% the duplicate pair is transcribed once. Physical p. 406 resumes at printed
% p. 384. Physical pp. 408--409 then repeat physical pp. 406--407
% (printed pp. 384--385); that duplicate pair is likewise transcribed once.
% Coverage: Section 13.2, Maximally Symmetric Spaces: Uniqueness; equations
%           (13.2.1)--(13.2.20), including all unnumbered displays.
% Figures/tables/footnotes: none.
% Status: source-reviewed and compile-clean.
% Uncertainties: none.
% MODERNIZATION: R denotes the curvature scalar. The curvature sign is the
% target sign, so K is the usual sectional-curvature constant:
% R=N(N-1)K, R_{\mu\nu}=(N-1)K g_{\mu\nu}, and the constant-curvature
% Riemann tensor has the target sign. Repeated coordinate derivatives are
% compact except where the Jacobian structure is pedagogically essential.
% LITERAL-OPERATOR AUDIT: every continuation operator in the curvature
% contractions and the power-series uniqueness proof was checked against the
% source; the contracted Bianchi identity retains its full tensor structure.

\subsection{Maximally Symmetric Spaces: Uniqueness}\label{sec:13.2}

We now show that the maximally symmetric spaces are uniquely specified by a
``curvature constant'' \(K\), and by the numbers of eigenvalues of the metric
that are positive or negative. That is, \emph{given two maximally symmetric
metrics with the same \(K\) and the same numbers of eigenvalues of each sign,
it will always be possible to find a coordinate transformation that carries
one metric into the other}. Armed with this theorem, we shall be able in the
next section to carry out an exhaustive study of maximally symmetric spaces by
simply constructing such metrics in one convenient coordinate system.

We showed in the last section that at any given point \(x\) in a maximally
symmetric space, we can find Killing vectors for which \(\xi_\lambda(x)\)
vanishes and for which \(\nabla_\kappa\xi_\lambda(x)\) is an arbitrary
antisymmetric matrix. It follows then that the coefficient of
\(\nabla_\kappa\xi_\lambda(x)\) in Eq.~\eqref{eq:13.1.12} must have a
vanishing antisymmetric part; that is,
\begin{align}
  &-R^\lambda{}_{\rho\sigma\nu}\delta_\mu^\kappa
  +R^\lambda{}_{\mu\sigma\nu}\delta_\rho^\kappa
  -R^\lambda{}_{\sigma\rho\mu}\delta_\nu^\kappa
  +R^\lambda{}_{\nu\rho\mu}\delta_\sigma^\kappa
  \notag\\
  &\qquad =
  -R^\kappa{}_{\rho\sigma\nu}\delta_\mu^\lambda
  +R^\kappa{}_{\mu\sigma\nu}\delta_\rho^\lambda
  -R^\kappa{}_{\sigma\rho\mu}\delta_\nu^\lambda
  +R^\kappa{}_{\nu\rho\mu}\delta_\sigma^\lambda.
  \tag{13.2.1}\label{eq:13.2.1}
\end{align}
We also showed that at any given point \(x\) in a maximally symmetric space,
there exist Killing vectors for which \(\xi_\lambda(x)\) takes any values
we like, so \eqref{eq:13.1.12} and \eqref{eq:13.2.1} require that
\begin{equation}
  \nabla_\nu R^\lambda{}_{\sigma\rho\mu}
  =
  \nabla_\sigma R^\lambda{}_{\nu\rho\mu}.
  \tag{13.2.2}\label{eq:13.2.2}
\end{equation}
We actually only need to use \eqref{eq:13.2.1}, because we have shown in the
last section that a space that is isotropic about every point, and hence
satisfies \eqref{eq:13.2.1}, must also be homogeneous, and hence must also
satisfy \eqref{eq:13.2.2}.

Our first step in the proof is to use Eq.~\eqref{eq:13.2.1} to derive a
formula for the curvature tensor. Contracting \(\kappa\) with \(\mu\)
yields
\begin{align*}
  -N R^\lambda{}_{\rho\sigma\nu}
  +R^\lambda{}_{\rho\sigma\nu}
  -R^\lambda{}_{\sigma\rho\nu}
  +R^\lambda{}_{\nu\rho\sigma}
  =
  -R^\lambda{}_{\rho\sigma\nu}
  +R_{\rho\sigma}\delta_\nu^\lambda
  -R_{\nu\rho}\delta_\sigma^\lambda.
\end{align*}
(Recall that \(R^\kappa{}_{\kappa\sigma\nu}\) vanishes,
\(R^\kappa{}_{\sigma\rho\kappa}=-R_{\sigma\rho}\), and in \(N\)
dimensions, \(\delta_\kappa^\kappa=N\).) Using the cyclic sum rule
\eqref{eq:6.6.5} and the antisymmetry of
\(R^\lambda{}_{\rho\sigma\nu}\), we find
\begin{equation}
  (N-1)R_{\lambda\rho\sigma\nu}
  =
  R_{\nu\rho}g_{\lambda\sigma}
  -R_{\sigma\rho}g_{\lambda\nu}.
  \tag{13.2.3}\label{eq:13.2.3}
\end{equation}
But this must be antisymmetric in \(\lambda\) and \(\rho\), so
\[
  R_{\nu\rho}g_{\lambda\sigma}
  -R_{\sigma\rho}g_{\lambda\nu}
  =
  -R_{\nu\lambda}g_{\rho\sigma}
  +R_{\sigma\lambda}g_{\rho\nu}.
\]
Contracting \(\lambda\) with \(\nu\), we find
\[
  R_{\sigma\rho}-N R_{\sigma\rho}
  =
  -R g_{\sigma\rho}+R_{\rho\sigma}.
\]
The Ricci tensor thus takes the form
\begin{equation}
  R_{\sigma\rho}
  =
  \frac{1}{N}g_{\sigma\rho}R.
  \tag{13.2.4}\label{eq:13.2.4}
\end{equation}
Inserting this in \eqref{eq:13.2.3} gives our formula for the curvature
tensor:
\begin{equation}
  R_{\lambda\rho\sigma\nu}
  =
  \frac{R}{N(N-1)}
  \left(
    g_{\nu\rho}g_{\lambda\sigma}
    -g_{\sigma\rho}g_{\lambda\nu}
  \right).
  \tag{13.2.5}\label{eq:13.2.5}
\end{equation}
This formula satisfies \eqref{eq:13.2.1}, so there is nothing further to be
learned from that condition.

In a space that is isotropic about \emph{every} point,
Eqs.~\eqref{eq:13.2.4} and \eqref{eq:13.2.5} will hold everywhere, and
we can use the Bianchi identities to say something about the dependence of
the curvature scalar \(R\) on position. Using \eqref{eq:13.2.4} in
\eqref{eq:6.8.4}, we have
\[
  0
  =
  \nabla_\rho
  \left(
    R^\rho{}_\sigma-\frac12\delta^\rho_\sigma R
  \right)
  =
  \left(\frac1N-\frac12\right)\partial_\sigma R,
\]
or
\begin{equation}
  0
  =
  \left(\frac1N-\frac12\right)\partial_\sigma R.
  \tag{13.2.6}\label{eq:13.2.6}
\end{equation}
Hence any space of three or more dimensions, in which
\eqref{eq:13.2.4} holds everywhere, will have \(R\) constant. It is
convenient to introduce a curvature constant \(K\) in place of \(R\), with
\begin{equation}
  R=N(N-1)K.
  \tag{13.2.7}\label{eq:13.2.7}
\end{equation}
Using this in \eqref{eq:13.2.4} gives the Ricci tensor and the
Riemann--Christoffel tensor here as
\begin{equation}
  R_{\sigma\rho}=(N-1)K g_{\sigma\rho},
  \tag{13.2.8}\label{eq:13.2.8}
\end{equation}
\begin{equation}
  R_{\lambda\rho\sigma\nu}
  =
  K\left(
    g_{\nu\rho}g_{\lambda\sigma}
    -g_{\sigma\rho}g_{\lambda\nu}
  \right).
  \tag{13.2.9}\label{eq:13.2.9}
\end{equation}
In differential geometry a space with these properties is called a
\emph{space of constant curvature}.

Incidentally, we showed in Section~\ref{sec:6.7} that the curvature tensor in
two dimensions is always of the form \eqref{eq:13.2.5}, so it is not
surprising that in this case \eqref{eq:13.2.6} does not allow us to draw any
conclusions about the constancy of \(R\). However, by using
\eqref{eq:13.2.2} one can show that the quantity \(K\) in
\eqref{eq:13.2.9} is also constant for maximally symmetric spaces of
dimensionality \(N=2\).

Now suppose that we are given two metrics \(g_{\mu\nu}(x)\) and
\(g'_{\mu\nu}(x')\), both having the same numbers of positive and negative
eigenvalues, and both satisfying the condition \eqref{eq:13.2.9} for a
maximally symmetric space; that is,
\begin{equation}
  R_{\lambda\rho\sigma\nu}
  =
  K\left(
    g_{\nu\rho}g_{\lambda\sigma}
    -g_{\sigma\rho}g_{\lambda\nu}
  \right),
  \tag{13.2.10}\label{eq:13.2.10}
\end{equation}
\begin{equation}
  R'_{\lambda\rho\sigma\nu}
  =
  K\left(
    g'_{\nu\rho}g'_{\lambda\sigma}
    -g'_{\sigma\rho}g'_{\lambda\nu}
  \right),
  \tag{13.2.11}\label{eq:13.2.11}
\end{equation}
with the same curvature constant \(K\). We shall show that
\(g_{\mu\nu}(x)\) and \(g'_{\mu\nu}(x')\) must be equivalent, in the sense
that there is a transformation \(x\longrightarrow x'\) that converts
\(g_{\mu\nu}(x)\) into \(g'_{\mu\nu}(x')\); that is, for which
\begin{equation}
  g'_{\mu\nu}(x')
  \frac{\partial x'^\mu}{\partial x^\rho}
  \frac{\partial x'^\nu}{\partial x^\sigma}
  =
  g_{\rho\sigma}(x).
  \tag{13.2.12}\label{eq:13.2.12}
\end{equation}
We shall prove this by actually constructing \(x'^\mu(x)\) as a power series
in \(x^\mu\). First, note that the equality in the numbers of positive and
negative eigenvalues of \(g_{\mu\nu}\) and \(g'_{\mu\nu}\) means that we
can find a nonsingular matrix \(d^\mu{}_\rho\) for which
\begin{equation}
  g'_{\mu\nu}(0)d^\mu{}_\rho d^\nu{}_\sigma
  =
  g_{\rho\sigma}(0).
  \tag{13.2.13}\label{eq:13.2.13}
\end{equation}
(The argument here is the same as in Section~\ref{sec:6.4}.) Thus we can
satisfy \eqref{eq:13.2.12} to zero order in \(x\) with
\[
  x'^\mu=d^\mu{}_\rho x^\rho.
\]

Now we proceed by mathematical induction. Suppose that we succeed in
satisfying \eqref{eq:13.2.12} to order \(n-1\) in \(x^\mu\) with a
polynomial
\begin{equation}
  x'^\mu(x)
  =
  d^\mu{}_\rho x^\rho
  +\sum_{m=2}^{n}\frac{1}{m!}
  d^\mu{}_{\rho_1\cdots\rho_m}
  x^{\rho_1}\cdots x^{\rho_m}.
  \tag{13.2.14}\label{eq:13.2.14}
\end{equation}
We want to add a term of order \(n+1\) in \(x^\mu\), so that
\eqref{eq:13.2.12} holds to order \(n\). This condition will be satisfied
if the derivative of \eqref{eq:13.2.12} holds in order \(n-1\); that is,
if
\begin{align*}
  &\frac{\partial^2x'^\mu}{\partial x^\rho\partial x^\lambda}
   \frac{\partial x'^\nu}{\partial x^\sigma}g'_{\mu\nu}(x')
  +
  \frac{\partial^2x'^\nu}{\partial x^\sigma\partial x^\lambda}
  \frac{\partial x'^\mu}{\partial x^\rho}g'_{\mu\nu}(x')
  \\
  &\qquad
  +
  \frac{\partial x'^\mu}{\partial x^\rho}
  \frac{\partial x'^\nu}{\partial x^\sigma}
  \frac{\partial x'^\kappa}{\partial x^\lambda}
  \frac{\partial g'_{\mu\nu}(x')}{\partial x'^\kappa}
  =
  \partial_\lambda g_{\rho\sigma}(x)
  \qquad\text{in order \(x^{n-1}\)}.
\end{align*}
This will be satisfied if (and, in fact, only if)
\begin{align*}
  \frac{\partial^2x'^\mu}{\partial x^\rho\partial x^\lambda}
  \frac{\partial x'^\nu}{\partial x^\sigma}g'_{\mu\nu}(x')
  &=
  g_{\sigma\tau}(x)\Gamma^\tau{}_{\lambda\rho}(x)
  \\
  &\quad
  -
  \frac{\partial x'^\mu}{\partial x^\rho}
  \frac{\partial x'^\nu}{\partial x^\sigma}
  \frac{\partial x'^\kappa}{\partial x^\lambda}
  g'_{\nu\tau}(x')\Gamma'^\tau{}_{\mu\kappa}(x')
  \qquad\text{in order \(x^{n-1}\)}.
\end{align*}
This only needs to hold in order \(n-1\) in \(x^\mu\), so we can use
\eqref{eq:13.2.12}, which was assumed to hold to this order, to convert it
into an equivalent requirement:
\begin{equation}
  \frac{\partial^2x'^\mu}{\partial x^\rho\partial x^\lambda}
  =
  \frac{\partial x'^\mu}{\partial x^\kappa}
  \Gamma^\kappa{}_{\lambda\rho}(x)
  -
  \frac{\partial x'^\nu}{\partial x^\rho}
  \frac{\partial x'^\kappa}{\partial x^\lambda}
  \Gamma'^\mu{}_{\nu\kappa}(x')
  \quad\text{in order \(x^{n-1}\)}.
  \tag{13.2.15}\label{eq:13.2.15}
\end{equation}
We can use \eqref{eq:13.2.14}, which is correct to order \(x^n\), to
calculate the term on the right-hand side of order \(x^{n-1}\). Let us write
the result as
\begin{align}
  &\left[
    \frac{\partial x'^\mu}{\partial x^\kappa}
    \Gamma^\kappa{}_{\lambda\rho}(x)
    -
    \frac{\partial x'^\nu}{\partial x^\rho}
    \frac{\partial x'^\kappa}{\partial x^\lambda}
    \Gamma'^\mu{}_{\nu\kappa}(x')
  \right]_{\text{order }n-1}
  \notag\\
  &\qquad =
  \frac{1}{(n-1)!}
  c^\mu{}_{\lambda\rho\sigma_1\cdots\sigma_{n-1}}
  x^{\sigma_1}\cdots x^{\sigma_{n-1}},
  \tag{13.2.16}\label{eq:13.2.16}
\end{align}
the coefficients
\(c^\mu{}_{\lambda\rho\sigma_1\cdots\sigma_{n-1}}\) depending in a
complicated way on the functions \(g_{\mu\nu}(x)\) and
\(g'_{\mu\nu}(x')\), and on the previously determined coefficients
\(d^\mu{}_{\rho_1\cdots\rho_m}\). Then \eqref{eq:13.2.15} will be
satisfied in order \(n-1\) if we add to \eqref{eq:13.2.14} a term
\begin{equation}
  \bigl[x'^\mu(x)\bigr]_{\text{order }n+1}
  =
  \frac{1}{(n+1)!}
  c^\mu{}_{\lambda\rho\sigma_1\cdots\sigma_{n-1}}
  x^\lambda x^\rho x^{\sigma_1}\cdots x^{\sigma_{n-1}},
  \tag{13.2.17}\label{eq:13.2.17}
\end{equation}
\emph{provided} that the coefficient
\(c^\mu{}_{\lambda\rho\sigma_1\cdots\sigma_{n-1}}\) is totally symmetric
in all its lower indices. These coefficients are obviously symmetric under
interchange of \(\lambda\) and \(\rho\), or among the \(\sigma_m\)
indices, so the only condition that needs to be satisfied is that they are
symmetric between \(\lambda\) and any \(\sigma_m\), or equivalently, that
the derivative of \eqref{eq:13.2.16} with respect to \(x^\sigma\) should
be symmetric between \(\lambda\) and \(\sigma\):
\begin{align}
  &\partial_\sigma
  \left(
    \frac{\partial x'^\mu}{\partial x^\kappa}
    \Gamma^\kappa{}_{\lambda\rho}(x)
    -
    \frac{\partial x'^\nu}{\partial x^\rho}
    \frac{\partial x'^\kappa}{\partial x^\lambda}
    \Gamma'^\mu{}_{\nu\kappa}(x')
  \right)
  \notag\\
  &\qquad =
  \partial_\lambda
  \left(
    \frac{\partial x'^\mu}{\partial x^\kappa}
    \Gamma^\kappa{}_{\sigma\rho}(x)
    -
    \frac{\partial x'^\nu}{\partial x^\rho}
    \frac{\partial x'^\kappa}{\partial x^\sigma}
    \Gamma'^\mu{}_{\nu\kappa}(x')
  \right)
  \qquad\text{in order \(x^{n-2}\)}.
  \tag{13.2.18}\label{eq:13.2.18}
\end{align}
Since \eqref{eq:13.2.12} is assumed to hold to order \(x^{n-1}\), its
derivative, Eq.~\eqref{eq:13.2.15}, will hold to order \(x^{n-2}\), so
we can use \eqref{eq:13.2.12} and \eqref{eq:13.2.15} to rewrite
\eqref{eq:13.2.18} as the equivalent requirement
\begin{equation}
  \frac{\partial x'^\mu}{\partial x^\kappa}
  R^\kappa{}_{\rho\lambda\eta}(x)
  =
  \frac{\partial x'^\nu}{\partial x^\rho}
  \frac{\partial x'^\kappa}{\partial x^\lambda}
  \frac{\partial x'^\sigma}{\partial x^\eta}
  R'^\mu{}_{\nu\kappa\sigma}(x')
  \qquad\text{in order \(x^{n-2}\)}.
  \tag{13.2.19}\label{eq:13.2.19}
\end{equation}
Now for the first time we use Eqs.~\eqref{eq:13.2.10} and
\eqref{eq:13.2.11}, which allow \eqref{eq:13.2.19} to be replaced with
the equivalent requirement
\begin{align}
  &\frac{\partial x'^\mu}{\partial x^\lambda}g_{\rho\eta}(x)
  -
  \frac{\partial x'^\mu}{\partial x^\eta}g_{\rho\lambda}(x)
  \notag\\
  &\qquad =
  \frac{\partial x'^\nu}{\partial x^\rho}
  \left(
    \frac{\partial x'^\mu}{\partial x^\lambda}
    \frac{\partial x'^\sigma}{\partial x^\eta}
    g'_{\nu\sigma}(x')
    -
    \frac{\partial x'^\kappa}{\partial x^\lambda}
    \frac{\partial x'^\mu}{\partial x^\eta}
    g'_{\nu\kappa}(x')
  \right)
  \qquad\text{in order \(x^{n-2}\)}.
  \tag{13.2.20}\label{eq:13.2.20}
\end{align}
This condition \emph{is} satisfied, because \eqref{eq:13.2.12} was
assumed to hold to order \(x^{n-1}\). To recapitulate, this implies that
\eqref{eq:13.2.19} holds in order \(x^{n-2}\), which implies that
\eqref{eq:13.2.18} holds in order \(x^{n-2}\), which implies that the
coefficients \(c^\mu{}_{\lambda\rho\sigma_1\cdots\sigma_{n-1}}\) are
totally symmetric in their lower indices, which implies that
\eqref{eq:13.2.17} satisfies \eqref{eq:13.2.15}, which implies that by
adding \eqref{eq:13.2.17} to \eqref{eq:13.2.14} we can satisfy
\eqref{eq:13.2.12} to order \(x^n\). Thus, if \eqref{eq:13.2.12} can
be satisfied to order \(x^{n-1}\) by a polynomial \(x'(x)\) of order
\(n\), it can be satisfied to order \(x^n\) by a polynomial \(x'(x)\) of
order \(n+1\), and therefore a function \(x'(x)\) satisfying
\eqref{eq:13.2.12} exactly can be built up as a power series, as was to be
proven.
